%!TEX root = ../thesis.tex

\chapter{Experimental Design and Reproducibility Methodology}
\label{chap:methodology}

\section{Test-case matrix}

Every comparison below is bound to its case, grid, solver, precision, final time
and baseline, and none transfers between unmatched configurations.
The ideal-MHD matrix pairs the circularly polarised Alfv\'en wave, which
carries the order verification, with the Brio--Wu shock tube and the periodic
Orszag--Tang and Kelvin--Helmholtz problems, covering a smooth exact solution,
a discontinuity, interacting waves and shear instability. Sod
\citep{toro2009} and Liska--Wendroff Configuration~3
\citep{liska_wendroff_2003} provide Euler continuity checks for the
cross-system comparison without repeating Report~1 validation, whose five-case
Euler coverage stands unchanged behind them. The Euler axes that report tested
at length, device, optimisation level, fast math and the HLLC branch rule, are
carried forward rather than rerun; the two cases kept here are those readable on
the same footing as their MHD counterparts.
Table~\ref{tab:ch3-case-matrix} sets out the completed coverage and case roles.

\begin{table}[htbp]
\centering
\small
\begin{tabularx}{\textwidth}{@{}lllclX@{}}
\toprule
Case & System & Dimension & Main grids & Solver & Role \\
\midrule
CP Alfv\'en & MHD & 1D & 32--512 & HLL$^{\ast}$ & exact solution, order verification \\
Sod & Euler & 1D & 200 & HLLC & continuity sensitivity \\
LW3 & Euler & 2D & $200^2$ & HLLC & continuity sensitivity \\
Brio--Wu & MHD & 1D & 200--8000 & HLL$^{\ast}$/HLLD & reference, device, precision, MCA \\
Orszag--Tang & MHD & 2D & $64^2$--$512^2$ & HLL$^{\ast}$/HLLD & refinement, device, precision, time \\
Kelvin--Helmholtz & MHD & 2D & $64^2$--$512^2$ & HLL$^{\ast}$/HLLD & refinement, precision, CFL, timing \\
\bottomrule
\end{tabularx}
\caption[Completed test-case matrix]{\textbf{Completed test-case matrix.}
The Euler rows provide continuity only. HLL$^{\ast}$ throughout is the globally
clamped path of Section~\ref{sec:ch2-riemann}, equivalent to a Rusanov flux
carrying $c_h$, and is written so wherever a table reports it as a solver
level.}
\label{tab:ch3-case-matrix}
\end{table}

\section{Build and run matrix}

State arrays, fluxes, wave-speed estimates and the CFL calculation follow the
build precision, while the time accumulator is binary64 in both builds, so that
$t \leftarrow t+\Delta t$ does not lose small increments in fp32; the fp32
build is a declared mixed-precision configuration, not a uniform binary32
program. The step is set by
\begin{equation}
 c_h=\max_{ij}\max\!\left(|v_x|+c_{f,x},\;|v_y|+c_{f,y}\right),\qquad
 \Delta t=\mathrm{CFL}\,\frac{\min(\Delta x,\Delta y)}{c_h},
 \label{eq:ch3-cfl}
\end{equation}
so one consequence carries through every comparison below: the stepwise
directional maximum that sets $c_h$ is formed in the build precision and then
sets $\Delta t$, so a precision or build change perturbs the time-step sequence
as well as the per-operation rounding, so a reported discrepancy is the response
to a declared axis rather than a pure per-cell rounding residual. Recorded step
counts nonetheless coincide across variants
(Section~\ref{sec:ch5-solver-timing}), so the paths differ in $\Delta t$ and not
in their number. Matched builds vary the harness-recorded optimisation,
floating-point-math and branch-rule semantics, since directory names do not
prove compiler behaviour. In MSVC terms O2 selects /O2, O3 /Ox and the composite
Ofast /Ox /fp:fast, which is why it confounds optimisation with math semantics;
the isolated matrix compares /O2 with /Ox, default math with /fp:fast, and
$\leq$ with $<$, one axis at a time. Virtual p53 and p24 MCA
perturb operations at 53 or 24 significand bits.

Table~\ref{tab:ch3-axis-matrix} records what each axis can support; these
metrics are not combined into one ranking.

\begin{table}[htbp]
\centering
\small
\begin{tabularx}{\textwidth}{@{}lllX@{}}
\toprule
Axis & Levels & Matched baseline & Reported quantity \\
\midrule
Precision & fp32/fp64 & same solver/build/grid & field discrepancy \\
Build & optimisation/math/branch & one axis changed & field discrepancy \\
Solver & HLL$^{\ast}$/HLLD & same case controls & timing or within-solver diagnostics \\
Device & CPU/GPU & same HLL$^{\ast}$ configuration & ULP, $L_\infty$, repeated timing \\
Resolution & 128/256/512 & adjacent same-group grid & block-averaged self-difference \\
Time & aligned checkpoints & same case and solver & fixed-window discrepancy fit \\
Robustness & requested threads/CFL & serial baseline/one setting changed & unused-setting null control or precision separation \\
MCA & p53/p24 & same stochastic scope & maximum cellwise spread \\
\bottomrule
\end{tabularx}
\caption[Controlled-axis comparison contract]{\textbf{Controlled-axis
comparison contract.}}
\label{tab:ch3-axis-matrix}
\end{table}

\section{Reference hierarchy and validation criteria}
\label{sec:ch3-reference-hierarchy}

Unit tests, transverse invariance and CPU/GPU agreement verify implementation
properties, not ideal MHD against physical reality
\citep{oberkampfRoy2010vv}. Evidence is ordered by independence: an exact
solution where one exists, then external fieldwise comparison,
high-resolution comparison, self-refinement, and property checks. Only the
circularly polarised Alfv\'en wave of Section~\ref{sec:ch4-cp-alfven} reaches
the first level, which is why it alone carries an order verification. Brio--Wu
uses an aligned fp64 $N=8000$ comparator, and block averaging aligns adjacent
two-dimensional grids; since those grids are not fine enough for a leading
truncation term to dominate, their fitted orders are descriptive rather than
Richardson or Grid Convergence Index estimates \citep{roache1994gci}.
A run enters measurement only on reaching the declared final time with a
positive step count, finite variables with $\rho,p>0$, and fresh output and
completion metadata.

Brio--Wu \citep{brio_wu_1988} uses outflow boundaries; the two-dimensional cases
below are doubly periodic with $\gamma=5/3$. The Brio--Wu data are
\begin{equation}
 (\rho,v_x,p,B_y)=
 \begin{cases}
  (1,\;0,\;1,\;1), & x<x_0,\\
  (0.125,\;0,\;0.1,\;-1), & x\ge x_0,
 \end{cases}
 \qquad
 \begin{aligned}
  &x\in[0,1],\; x_0=0.5,\; \gamma=2,\; t=0.1,\\
  &B_x=0.75,\; v_y=v_z=B_z=\psi=0 .
 \end{aligned}
 \label{eq:ch3-ic-briowu}
\end{equation}
For Orszag--Tang at $t=0.5$,
\begin{equation}
 \begin{aligned}
  \rho &= \gamma^{2}=25/9,\qquad p=\gamma=5/3,\\
  \mathbf v &= (-\sin 2\pi y,\;\sin 2\pi x),\\
  \mathbf B &= B_0(-\sin 2\pi y,\;\sin 4\pi x),\qquad B_0=1,\\
  &v_z=B_z=\psi=0,\qquad c_r=0.18,\qquad
  (x,y)\in[0,1]^2\ \text{(periodic)} .
 \end{aligned}
 \label{eq:ch3-ic-ot}
\end{equation}
Under $X=2\pi x$, $Y=2\pi y$, $T=2\pi t$ this is the original $[0,2\pi]^2$
benchmark at $T=\pi$ \citep{orszagTang1979} in the normalisation of
\citet{toth_2000}, and its separated dependencies $B_x(y)$, $B_y(x)$ give
exactly zero cell-centred divergence initially, so later values measure
numerical divergence. The project-defined smooth Kelvin--Helmholtz double shear
layer at $t=1.0$ is
\begin{align}
 &(x,y)\in[0,1]^2\ \text{(periodic)},\quad
 \rho=p=1,\quad \gamma=5/3,\quad c_r=0.18, \notag\\
 v_x &= U_0\left[\tanh\frac{y-y_1}{a}-\tanh\frac{y-y_2}{a}-1\right], \notag\\
 v_y &= \delta\sin(2\pi x)\left[\mathrm{e}^{-(y-y_1)^2/s^2}\right. \notag\\
     &\hspace{5.6em}\left.{}+\mathrm{e}^{-(y-y_2)^2/s^2}\right], \notag\\
 &(y_1,y_2)=(0.25,0.75),\quad U_0=0.5,\quad a=0.025, \notag\\
 &\delta=0.01,\quad s=0.05,\quad B_x=B_0=0.1, \notag\\
 &v_z=B_y=B_z=\psi=0 .
 \label{eq:ch3-ic-kh}
\end{align}
The field is uniform and flow-aligned at Alfv\'en Mach number $5$, so the
initial state is exactly divergence-free.

The refinement ladder uses $\mathrm{CFL}=0.4$ for HLL and $0.2$ for HLLD
(Section~\ref{sec:ch4-orszag-tang}); HLL was also re-run at $0.2$ so the solver
axis can be read at a matched time step.

\section{Metrics and statistical treatment}
\label{sec:ch3-metrics}

For a same-grid difference $d_{ij}=q_{ij}-q^{\mathrm{ref}}_{ij}$ over
$N=N_xN_y$ cells, the reported mean norms are
\begin{equation}
 \begin{aligned}
 L_{1,\mathrm{mean}}&=\frac{1}{N}\sum_{ij}|d_{ij}|,\\
 L_{2,\mathrm{mean}}&=\left(\frac{1}{N}\sum_{ij}d_{ij}^{2}\right)^{1/2},
 \end{aligned}
 \label{eq:ch3-mean-norms}
\end{equation}
with $L_\infty=\max_{ij}|d_{ij}|$. The compact cross-system matrix uses the
mean-relative density difference
\begin{equation}
 L_{1,\mathrm{rel}}(\rho)=
 \frac{N^{-1}\sum_{ij}|\rho_{ij}-\rho^{\mathrm{ref}}_{ij}|}
      {N^{-1}\sum_{ij}|\rho^{\mathrm{ref}}_{ij}|},
 \label{eq:ch3-mean-relative}
\end{equation}
and only within that matrix. For binary32 the unit roundoff is
$u_{32}=2^{-24}\approx5.96\times10^{-8}$ \citep{ieee754_2019}; discrepancies are
also reported as $L_{1,\mathrm{rel}}/u_{32}$, so a value near unity means the
two calculations differ by about one rounding of the field, and that ratio is
written below as a number of roundings. Since $L_{1,\mathrm{rel}}$ is normalised
by the reference field mean whereas $L_\infty$ is absolute, the two are not on a
common scale. Adjacent-grid comparisons use block-averaged
absolute density $L_{1,\mathrm{mean}}$ with the descriptive diagnostic
$p_{\mathrm{obs}}=\log_2(E_{128,256}/E_{256,512})$. Writing
$D_N=\|\rho_{32,N}-\rho_{64,N}\|_{1,\mathrm{mean}}$ and $E^{64}_{256,512}$ for
the matched fp64 adjacent-grid disagreement, the ratio
$S_N=D_N/E^{64}_{256,512}$ measures neither exact error, nor GCI uncertainty,
nor fp32 adequacy.

The divergence diagnostic applies centred differences on interior cells,
\begin{align}
 D_{ij} &=\frac{B_{x,i+1,j}-B_{x,i-1,j}}{2\Delta x}
          +\frac{B_{y,i,j+1}-B_{y,i,j-1}}{2\Delta y}, \notag\\
 \mathrm{divB}_{\mathrm{mean}}
        &=\frac{1}{N_I}\sum_{(i,j)\in I}|D_{ij}|,
 &\mathrm{divB}_{\max}
        &=\max_{(i,j)\in I}|D_{ij}|.
 \label{eq:ch3-divb}
\end{align}
where $I$ excludes the outer stencil layer. These carry dimensions of magnetic
field per length. The $2\Delta x$ stencil is not the compact operator the GLM
source controls and is blind to grid-scale oscillation in $\mathbf B$, so small
values bound one discrete residual rather than field accuracy
\citep{toth_2000}. Same-precision CPU/GPU comparisons use absolute $L_\infty$
and the maximum distance in units in the last place (ULP) on saved conservative
states, mapping sign-adjusted bit patterns to monotone integers. Zero ULP
denotes bitwise equality; ULP is neither a cross-precision nor a physical error
norm.

For $n$ MCA samples the per-cell sample standard deviation is
\begin{equation}
 \sigma_j=\left[\frac{1}{n-1}\sum_{s=1}^{n}
 (q_{s,j}-\bar q_j)^2\right]^{1/2},\qquad
 \mathrm{spread}_q=\max_j\sigma_j ,
 \label{eq:ch3-mca-spread}
\end{equation}
the maximum conservatively aggregating shocks and shear layers.

Elapsed time is subprocess wall time from launch through completion and
required output, not kernel time. The KH solver/precision groups discard one
CPU warm-up and retain $n=5$; the CPU/GPU packet retains $n=5$ but excludes no
device warm-up or outlier, so context creation is inside its Brio--Wu figures.
Both report median and interquartile range, for the recorded workstation only.

The temporal experiment aligns fp32 and fp64 saved states at each time and
computes density $L_{1,\mathrm{mean}}$ and $L_\infty$. For positive samples in
the predeclared case-specific window, ordinary least squares fits two
competing models,
\begin{equation}
 \log e(t)=a+\lambda t,
 \qquad
 \log e(t)=a'+k\log t .
 \label{eq:ch3-temporal-fit}
\end{equation}
Here $t$ is code time, so $1/\lambda$ is the e-folding time. The exponent $k$
instead reads the series as accumulation: at fixed CFL the step count is
proportional to $t$, so coherent per-step addition gives $k\simeq1$ and an
incoherent random walk $k\simeq0.5$ \citep{crociEtAl2022stochastic}, whereas
exponential growth is the signature of amplification. Both models carry two free
parameters and are fitted to identical samples, so the smaller root-mean-square
log residual identifies the better-supported one without an information
criterion; a near-zero $R^2$ records only the absence of a slope and
cannot separate a flat series from a noisy one, whereas that residual separates
a plateau from scatter. Fitted exponents carry the ordinary least-squares
standard error of the slope, which is optimistic because successive stopping
times of one case are not independent samples, and both $\lambda$ and $k$ stay
descriptive rates, not physical instability rates.

\section{Harness, metadata, and retention}
\label{sec:ch3-harness-metadata}

Experiments follow configuration, build, run, measurement, aggregation and
plotting. Each record fixes the commit, the configuration and binary hashes, the
compiler and effective floating-point semantics, the environment, and the return
status, completion diagnostics and wall time that separate a complete run from a
truncated one. Aggregation checks matrix completeness and physical admissibility
before any summary or figure. Retention is listed in
Appendix~\ref{sec:appendix-provenance}.
